\chapter{The game is a market}\label{ch:market}
\providecommand{\Xlo}{\underline{X}}\providecommand{\Xhi}{\overline{X}}

\paragraph{Idea.} In an equilibrium the coverage plays the role of prices and the attack loads the role
of demand (Section~\ref{sec:intro-example}). This chapter makes the analogy exact. In one direction, a
market can be hidden inside a game with a single resource, using extra ``private'' and ``control''
targets that pin down the defender's behavior. In the other direction, a game becomes a market once the
defender's cutoff $\lambda$ is known (Lemma~\ref{lem:cutoff}): each target $j$ is then a good with
capacity $\lambda/d_j$. The cutoff is not known, so we search for it by bisection, and since the market
oracle may return any of many equilibria, we need a lattice structure, a ``pinning'' trick and a final
linear program to hit the budget $\sum_jx_j=r$ exactly.

\paragraph{Example.} In the example of Figure~\ref{fig:example-market} the defender gains are $d_j=1$
and the cutoff is $\lambda=1/2$. The equilibrium is then exactly an equilibrium of the market in which
the two types are buyers of masses $2/3$ and $1/3$, the two targets are goods of capacity
$\lambda/d_j=1/2$, and the prices are $x=(3/4,1/4)$: both prices lie strictly inside $[0,1]$, so both
goods must sell exactly their capacity, and they do. The budget condition $x_1+x_2=1$ is what singles
out the cutoff $1/2$; finding that cutoff is the main work of Section~\ref{sec:mk-g2m}.

\section{The main theorem}

\begin{theorem}[Game $\Leftrightarrow$ market]\label{thm:mk-main}
\statusrefereed
Exact BNE search for explicit Bayesian simple security games and exact equilibrium search for
explicit (M1) markets with arbitrary positive rational slopes are polynomial-time Turing equivalent.
More precisely:
\begin{enumerate}[nosep,label=(\alph*)]
\item \emph{Market to game.} From an (M1) instance one computes in polynomial time a game with
\emph{one} defender resource and defender payoffs $u_1^0\equiv0$, $u_1^1\equiv1$, such that \emph{every}
BNE of it decodes in polynomial time to an equilibrium of the market. One call to a BNE oracle
suffices.
\item \emph{Game to market.} There is a polynomial-time algorithm that, with access to an oracle
returning \emph{any} equilibrium of a dense (M1) instance (of polynomial length, as above), computes an
exact BNE of an arbitrary explicit game, including games with zero slopes $b_{ij}=0$ and zero gains
$d_j=0$. It makes polynomially many oracle calls, each of polynomial binary length, followed by
rational linear programming.
\end{enumerate}
\end{theorem}

A polynomial-time algorithm for either problem therefore gives one for the other. Composing (b) with
(a) gives the following.

\begin{corollary}[One resource suffices]\label{cor:mk-one-resource}
\statusrefereed
Exact BNE search in general explicit Bayesian simple security games reduces, by a polynomial-time
Turing reduction, to exact BNE search in games with one defender resource, defender payoffs $0$
uncovered and $1$ covered at every target, and, for any fixed rational $0<\omega\le1$, type
probabilities satisfying $1-\omega\le|\Theta|\Pr(\theta)\le1+\omega$ (where $|\Theta|$ is the number of types)
with largest-to-smallest ratio at most $1+\omega$. The two problems are therefore polynomial-time
Turing equivalent.
\end{corollary}

So a polynomial algorithm for one-resource games with identical defender payoffs and nearly uniform
priors would already answer Conitzer's question. Section~\ref{sec:mk-g2m} proves (b),
Section~\ref{sec:mk-m2g} proves (a), and Section~\ref{sec:mk-proof} assembles the theorem and the
corollary. Chapter~\ref{ch:arrival} uses only Section~\ref{sec:mk-m2g}; Section~\ref{sec:mk-g2m}, the longer direction, is needed only for part~(b) and Corollary~\ref{cor:mk-one-resource}, and a first reading can skip it. 

\section{Game to market}\label{sec:mk-g2m}

\paragraph{Idea.} The proof has four steps.
\begin{enumerate}[nosep]
\item Fix a cutoff $\lambda$. The game's equilibrium conditions become those of a \emph{box market}
(Definition~\ref{def:mk-box}), and a box market is one (M1) call (Lemma~\ref{lem:mk-box-m1}).
\item As $\lambda$ grows the equilibrium coverage falls, and at a fixed $\lambda$ the equilibrium prices
form a lattice (Proposition~\ref{prop:mk-lattice}). Bisection on $\lambda$ then brackets the budget $r$
between two ordered equilibria $\Xlo{}\le \Xhi{}$ at one exact cutoff $\lambda_0$ (Lemma~\ref{lem:mk-bisect}).
\item Pinning one coordinate between $\Xlo{}$ and $\Xhi{}$ and solving the smaller box market gives a new
equilibrium in between (Lemma~\ref{lem:mk-pin}). Repeating this shrinks $[\Xlo{},\Xhi{}]$ until a linear program
can repair the budget exactly (Lemma~\ref{lem:mk-narrow}).
\item Zero slopes and gains are removed by a tiny perturbation (Lemma~\ref{lem:mk-degenerate}).
\end{enumerate}
Throughout, $p_i>0$ and $0<r<n$ (Section~\ref{sec:pre-game}).
Figure~\ref{fig:g2m-flow} shows how the pieces fit together.

\begin{figure}[t]
\centering
\begin{tikzpicture}[
  st/.style={draw,rounded corners,align=center,font=\small,minimum height=0.95cm,text width=3.7cm},
  or/.style={draw,thick,fill=gray!10,align=center,font=\small,minimum height=0.95cm,text width=3.0cm},
  ar/.style={-{Stealth[length=2mm]},thick},
  lb/.style={font=\scriptsize,fill=white,inner sep=1pt,align=center}]
  \node[st] (g) at (0,0) {explicit game; zero slopes/gains perturbed (Lemma~\ref{lem:mk-degenerate})};
  \node[st] (b) at (0,-1.9) {cutoff bisection on $\lambda$ (Lemma~\ref{lem:mk-bisect})};
  \node[st] (n) at (0,-3.8) {narrowing: pin one coordinate in $[\Xlo,\Xhi]$, canonicalize, patch (Lemmas~\ref{lem:mk-pin}, \ref{lem:mk-narrow})};
  \node[st] (r) at (0,-5.9) {exact budget repair: price LP of the pattern of $\Xlo$ (Lemma~\ref{lem:mk-narrow})};
  \node[st] (e) at (0,-7.8) {exact BNE (Lemma~\ref{lem:mk-box-cutoff}); if perturbed, decode by one LP (Lemma~\ref{lem:mk-degenerate}); realize by Lemma~\ref{lem:realize}};
  \node[or] (q) at (6.9,-2.85) {box-market query (Def.~\ref{def:mk-box}) $\to$ one dense (M1) call (Lemma~\ref{lem:mk-box-m1})};
  \draw[ar] (g) -- (b);
  \draw[ar] (b) -- node[lb,right]{bracket $\Xlo\le\Xhi$ at $\lambda_0$} (n);
  \draw[ar] (n) -- node[lb,right]{all widths $\le\eta$} (r);
  \draw[ar] (r) -- (e);
  \draw[{Stealth[length=2mm]}-{Stealth[length=2mm]},thick] (b.east) -| node[lb,above,pos=0.3]{midpoint query / any equilibrium} (q.north);
  \draw[{Stealth[length=2mm]}-{Stealth[length=2mm]},thick] (n.east) -| node[lb,below,pos=0.3]{pinned sub-box / any equilibrium} (q.south);
  \draw[ar,dashed] (b.west) to[out=180,in=180] node[lb,left]{total $=r$:\\stop} (e.west);
  \draw[ar,dashed] (n.west) to[out=200,in=160] (e.170);
  \node[font=\scriptsize,align=center,anchor=north] at (6.9,-4.5) {polynomially many calls, each of\\polynomial length (Lemma~\ref{lem:mk-bits})};
\end{tikzpicture}
\caption{The game-to-market reduction of Proposition~\ref{prop:mk-g2m}. Only the shaded box uses the (M1)
oracle; every other step is rational linear programming. Bisection and narrowing each query the oracle
repeatedly, and the oracle may return any equilibrium. Either loop may stop early with an exact BNE when an
answer has total exactly $r$ (dashed arrow).}
\label{fig:g2m-flow}
\end{figure}

\subsection{Box markets}

\begin{definition}[Box market]\label{def:mk-box}
Fix the types as buyers (masses $p_i$, payoffs $U_{ij}$), capacities $c_j\ge0$, and a price box
$l\le h$ in $\mathbb R^n$. An \emph{equilibrium of the box market} is a pair $(x,f)$ with $l\le x\le h$,
$f\ge0$, $\sum_jf_{ij}=p_i$ for all $i$, $f_{ij}>0$ only if $j$ maximizes $U_{ij}(x)$ over all of $[n]$,
and
\[
x_j>l_j\ \Longrightarrow\ w_j\ge c_j,\qquad x_j<h_j\ \Longrightarrow\ w_j\le c_j .
\]
A pinned coordinate $l_j=h_j$ imposes no condition on its load.
\end{definition}

\begin{lemma}[The game at a fixed cutoff]\label{lem:mk-box-cutoff}
Let $0<r<n$ and $d_j>0$ for all $j$. A pair $(x,f)$ is a BNE if and only if $\sum_jx_j=r$ and, for
some $\lambda\in[0,\max_jd_j]$, $(x,f)$ is an equilibrium of the box market on $[0,1]^n$ with
capacities $c_j=\lambda/d_j$.
\end{lemma}

\begin{proof}
Since $d_j>0$, $d_jw_j\ge\lambda$ if and only if $w_j\ge\lambda/d_j$, and likewise for $\le$. Apply
Lemma~\ref{lem:cutoff}.
\end{proof}

Until Section~\ref{sec:mk-degenerate} we assume $b_{ij}>0$ and $d_j>0$.

\subsection{Demand comparison and the price lattice}

\paragraph{Idea.} If the prices of a set $S$ of goods go up while all other prices weakly go down, the
demand for $S$ cannot rise. This simple monotonicity gives order between cutoffs and a lattice at a
fixed cutoff.

For a set $S$ of goods and a flow $f$ write $w_f(S)=\sum_{j\in S}\sum_if_{ij}$. All slopes are
strictly positive in this subsection.

\begin{lemma}[Demand comparison]\label{lem:mk-demand}
Suppose $x_j>y_j$ for $j\in S$ and $x_j\le y_j$ for $j\notin S$, and let $f$, $g$ be best-response flows
(full rows, positive only on row maximizers) at $x$ and $y$. Then $w_f(S)\le w_g(S)$. If equality holds,
exactly the same set of types sends mass to $S$ under $f$ and $g$, and each such type sends its entire
mass to $S$ under both.
\end{lemma}

\begin{proof}
If type $i$ sends mass to $S$ under $f$, some $j\in S$ has $U_{ij}(x)=V_i(x)$. Then
$U_{ij}(y)>U_{ij}(x)=V_i(x)\ge U_{ik}(x)\ge U_{ik}(y)$ for every $k\notin S$, so all best goods of $i$
at $y$ lie in $S$ and $i$ contributes $p_i$ to $w_g(S)$, at least its contribution to $w_f(S)$. Types
sending nothing to $S$ under $f$ contribute $0$ to $w_f(S)$. Summing over types with positive priors
gives the inequality. At equality every type's two contributions agree, which is the isolation
statement. The case $S=\emptyset$ is trivial.
\end{proof}

\begin{proposition}[Order across cutoffs and lattice at one cutoff]\label{prop:mk-lattice}
\statusrefereed
Consider box markets on $[0,1]^n$ with capacities $c_j=\lambda/d_j$.
\begin{enumerate}[nosep,label=(\alph*)]
\item If $(x,f)$ is an equilibrium at cutoff $\lambda'$ and $(y,g)$ one at cutoff $\lambda<\lambda'$,
then $x\le y$ coordinatewise. This holds for every pair of selections.
\item At fixed capacities $c$ (on any common box), if $(x,f)$ and $(y,g)$ are equilibria, then so are
$\min(x,y)$ and $\max(x,y)$, with explicitly patched flows: with $S=\{j:x_j>y_j\}$ and $C$ its
complement, use $g$ on $S$ and $f$ on $C$ for the minimum, and $f$ on $S$ and $g$ on $C$ for the maximum.
\end{enumerate}
\end{proposition}

\paragraph{Proof roadmap.} (a) applies Lemma~\ref{lem:mk-demand} to the set $S$ where $x>y$: higher prices on
$S$ cannot come with higher demand there. (b) has three steps: (1) Lemma~\ref{lem:mk-demand} forces equal
loads on $S$ and isolates the types serving $S$; (2) for $\min(x,y)$, take $g$ on $S$ and $f$ on its
complement and check every row through $V_i(\min(x,y))=\max(V_i(x),V_i(y))$; (3) the maximum is
symmetric.

\begin{proof}
(a) Let $S=\{j:x_j>y_j\}$. For $j\in S$, $x_j>0$ and $y_j<1$, so
$w_{f,j}\ge\lambda'/d_j>\lambda/d_j\ge w_{g,j}$. If $S\ne\emptyset$, summing over $S$ contradicts
Lemma~\ref{lem:mk-demand}.

(b) On $S$, $x_j>y_j$ gives $x_j>l_j$ and $y_j<h_j$, so $w_{f,j}\ge c_j\ge w_{g,j}$.
Lemma~\ref{lem:mk-demand} forces equality on every good of $S$ and isolates the same set $I$ of types on
$S$ under both flows. Let $z=\min(x,y)$, and let $q$ equal $g$ on the columns of $S$ and $f$ on the
columns of $C$. Since prices enter separately and payoffs decrease, $U_{ij}(z)=\max(U_{ij}(x),U_{ij}(y))$
and hence $V_i(z)=\max(V_i(x),V_i(y))$. For $i\in I$, a good of $S$ used at $x$ strictly improves at $y$,
so $V_i(y)>V_i(x)$; its $q$-goods are $g$-goods in $S$ with value $V_i(y)=V_i(z)$. For $i\notin I$, a good
it uses at $y$ lies in $C$, where $x\le y$, so $V_i(x)\ge V_i(y)$; its $q$-goods are $f$-goods in $C$ with
value $V_i(x)=V_i(z)$. Row sums hold by isolation. Every column takes its flow and its price from the
same source equilibrium, so its box condition holds. For $z'=\max(x,y)$, use $f$ on $S$ and $g$ on $C$:
a selected $f$-good in $S$ keeps its price from $x$ while all competitors can only become more
expensive, and the same holds for selected $g$-goods in $C$ relative to $y$. Rows and columns are as
before.
\end{proof}

For unit-mass one-to-one assignment with general decreasing utilities, the price lattice and continua
of equilibria have precedents in Demange and Gale~\cite{dg85} (the lattice) and in~\cite[\S4]{alaei12}.

\subsection{Pinning}

\paragraph{Idea.} Given two ordered equilibria $\Xlo{}\le \Xhi{}$, fix one coordinate $k$ at any value $t$ in
between and solve the box market on $[\Xlo{},\Xhi{}]$. Whatever the oracle returns can be patched, column by
column, with $\Xlo{}$ or $\Xhi{}$ into an equilibrium of the full box.

\begin{lemma}[Pinning]\label{lem:mk-pin}
Fix capacities $c$. Let $(\Xlo{},f^L)$ and $(\Xhi{},f^H)$ be equilibria of the box market on $[0,1]^n$ with
$\Xlo{}\le \Xhi{}$. Let $k$ be a coordinate and $\Xlo_k\le t\le \Xhi_k$. Let $(x,f)$ be \emph{any} equilibrium of the
box market on $[\Xlo{},\Xhi{}]$ with coordinate $k$ pinned to $t$ (that is, $l_k=h_k=t$). Then $x$ is an
equilibrium price of the box market on $[0,1]^n$, with a supporting flow that takes each column from
$f$, $f^L$ or $f^H$.
\end{lemma}

\paragraph{Proof roadmap.} Two cases, according to whether the pinned coordinate $k$ is over- or
under-demanded. In each case: (1) let $S$ be the set where $x$ differs from the bracket endpoint
($\Xlo$ or $\Xhi$); (2) Lemma~\ref{lem:mk-demand} forces load exactly $c_j$ on $S$ and isolates the
types serving $S$; (3) patch, taking columns of $S$ from $f$ and the other columns from the endpoint, and
check rows and columns.

\begin{proof}
\emph{Case $w_{f,k}\ge c_k$.} Let $S=\{j:x_j>\Xlo_j\}$. For $j\in S\setminus\{k\}$ the sub-box gives
$w_{f,j}\ge c_j$, and the hypothesis gives it at $k$. Since $\Xlo_j<x_j\le1$, the equilibrium $\Xlo{}$ has
$w_{f^L,j}\le c_j$ on $S$. Outside $S$, $x_j=\Xlo_j$. Lemma~\ref{lem:mk-demand} applied to $x$ and $\Xlo{}$
forces $w_{f,j}=c_j=w_{f^L,j}$ on $S$ and isolates the same set $I$ of types on $S$ under $f$ and $f^L$.
Use $f$ on the columns of $S$ and $f^L$ elsewhere. Types in $I$ use $f$ only, so their choices are best
at $x$. A type outside $I$ uses $f^L$-goods outside $S$, whose prices are unchanged from $\Xlo{}$, while all
other prices have only risen; so its choices remain best at $x$. Rows are full by isolation. Columns of
$S$ carry exactly $c_j$, which is admissible at every status in $[0,1]^n$, and columns outside $S$ keep
the price and load of $\Xlo{}$.

\emph{Case $w_{f,k}\le c_k$.} Let $S=\{j:x_j<\Xhi_j\}$. Here $w_{f,j}\le c_j$ on $S$ (by the sub-box for
$j\ne k$ and the hypothesis at $k$) and $w_{f^H,j}\ge c_j$ on $S$ since $\Xhi_j>x_j\ge0$. Outside $S$,
$x_j=\Xhi_j$. Lemma~\ref{lem:mk-demand} applied to $\Xhi{}$ and $x$ forces equality on $S$ and isolates the
types serving $S$. Use $f$ on $S$ and $f^H$ elsewhere. Types serving $S$ use $f$ only. A type not serving
$S$ under $f$ puts all its $f$-mass outside $S$, so its best value at $x$ is attained outside $S$; the
prices outside $S$ are the same at $\Xhi{}$ and $x$, so every $f^H$-good it uses outside $S$ has value
$V_i(\Xhi{})\ge\max_{k'\notin S}U_{ik'}(\Xhi{})=V_i(x)$, hence is best at $x$. Rows and columns are as in the first
case.
\end{proof}

In either case a supporting flow can also be recovered by a transportation feasibility LP at the
known prices $x$.

\begin{example}[Interpolation does not work]\label{ex:mk-interp}
One might hope to obtain the budget $r$ by interpolating between two ordered equilibria with totals
on both sides of $r$. Take three targets, one resource, priors $(4,2,1)/7$, gains $(1,2/3,1)$ and attack
payoffs
\[
\begin{pmatrix}
4/3-2x_1&4/3-x_2&-2-x_3\\
1-x_1&7/6-x_2&4/3-x_3\\
-2-x_1&-2-x_2&4/3-x_3
\end{pmatrix}.
\]
At $\lambda=2/7$, $\Xlo{}=(0,0,1/6)$ and $\Xhi{}=(1/3,2/3,2/3)$ are box equilibria, with attacks (in units of
$1/7$) $((2,2,0),(0,1,1),(0,0,1))$ and $((1,3,0),(1,0,1),(0,0,1))$, both with loads $(2,3,2)/7$ and all
scores $2/7$. Their totals are $1/6$ and $5/3$. The point $\tfrac49\Xlo{}+\tfrac59\Xhi{}=(5/27,10/27,4/9)$ has
total one, but no BNE uses it: there the second type's payoffs are $(22/27,43/54,8/9)$, so types two and
three both attack target three, whose score is then at least $3/7$, while all three coverages are
interior and so all scores would have to equal $1/(1+3/2+1)=2/7$. A BNE off the chord is
$x=(1/6,1/3,1/2)$ with the attacks of $\Xlo{}$, at which the second type is tied at $5/6$ on all targets. The
procedure below never interpolates; it makes further pinned queries and then solves an LP on the
support of an endpoint, which may leave the chord.
\end{example}

\subsection{Box queries are ordinary (M1) calls}

\begin{lemma}[Box query to (M1)]\label{lem:mk-box-m1}
Let $c_j>0$ for all $j$ and let $l\le h$ be rational. Form the dense (M1) instance in general form with
buyers $i$ of mass $p_i$, goods $j$ of capacity $c_j$, payoffs $a_{ij}-b_{ij}P_j$, seller reserves $l_j$
and buyer reserves
\[
v_i^0=\max_k\bigl(a_{ik}-b_{ik}h_k\bigr).
\]
From \emph{any} equilibrium $(P,v,f)$ of it one obtains in polynomial time an equilibrium of the box
market on $[l,h]$ with capacities $c$, with $x=P$.
\end{lemma}

\begin{proof}
Suppose $P_j>h_j$. Every buyer has $a_{ij}-b_{ij}P_j<a_{ij}-b_{ij}h_j\le v_i^0\le v_i$, so column $j$ is
empty; but $P_j>h_j\ge l_j$ forces it to be full with $c_j>0$. Hence $l\le P\le h$. A buyer with spare
mass has $v_i=v_i^0$. For a maximizer $k$ in its reserve formula, the no-blocking inequality gives
$a_{ik}-b_{ik}h_k=v_i\ge a_{ik}-b_{ik}P_k$, so $P_k\ge h_k$ and hence $P_k=h_k$; and $k$ is a best good of
$i$. Assign the spare mass to such a $k$. All used goods now maximize $U_{ij}(P)$, since $v_i$ is at
least every payoff and equal on used goods. A good below its ceiling received no extra mass and has
load at most $c_j$. A good above its floor had a full column and so still has load at least $c_j$.
Pinned coordinates have no load condition.
\end{proof}

Subtracting the reserves (Section~\ref{sec:pre-market}) turns this call into a dense (M1) instance in
the zero-reserve form, with intercepts $a_{ij}-b_{ij}l_j-v_i^0$, the original positive slopes, masses
$p_i$ and capacities $c_j$. Pinned goods and ceilings are encoded in the data through $v_i^0$ and the
floor shift; they are not part of the model. The capacities $c_j=\lambda/d_j$ are rationals written in
binary but in general not dyadic. If an oracle requires dyadic data, multiplying all masses and
capacities by a common integer (which preserves prices and utilities) is not enough by itself: one must
also clear the denominators of each buyer's intercepts and slopes by a row rescaling with a positive integer $\nu_i$ (the proof of Lemma~\ref{lem:mk-rescale} uses only $\nu_i>0$).

\subsection{Locating the cutoff exactly}

\paragraph{Idea.} Bisection on $\lambda$ needs only finitely many steps, because the oracle's answers
can be summarized by finitely many \emph{patterns}, and the cutoffs at which a pattern stops being
valid are rationals of bounded denominator.

Put $D=1+\max_jd_j$. At cutoff $0$ the corner $x=\mathbf1$ with any best responses is a box
equilibrium (all capacities are $0$); at cutoff $D$ every $c_j=D/d_j>1\ge w_j$, so $x=\mathbf0$ with any
best responses is one.

A \emph{pattern} of a box equilibrium $(x,f)$ at cutoff $\lambda$ is its set $A$ of positive edges
together with the label \emph{zero}, \emph{one} or \emph{free} of each target ($x_j=0$, $x_j=1$, or
$0<x_j<1$). The \emph{pattern polytope} of a pattern, with a sign $\sigma\in\{\ge,\le\}$, is the set of
$(x,f,\lambda)$ with $x\in[0,1]^n$, $f\in[0,1]^{m\times n}$, $0\le\lambda\le D$, $f_{ij}=0$ for
$(i,j)\notin A$, $\sum_jf_{ij}=p_i$, $U_{ij}(x)\ge U_{ik}(x)$ for $(i,j)\in A$ and all $k$; $x_j=0$ and
$d_jw_j\le\lambda$ at zero targets; $x_j=1$ and $d_jw_j\ge\lambda$ at one targets; $d_jw_j=\lambda$ at free
targets; and $\sum_jx_j\ \sigma\ r$. Zero flow on an edge of $A$ is allowed. Every point of it is a box
equilibrium at its own cutoff $\lambda$ with the indicated total inequality.

Clear the denominators of every row of the polynomial-size menu of all rows that can occur in pattern
polytopes, and let $H_1\ge1$ bound all absolute integer coefficients and right-hand sides. With
$N_1=n+mn+1$ variables put $Q_1=(N_1H_1)^{N_1}$; its binary length is polynomial in the input.

\begin{lemma}[Cutoff bisection with pattern intervals]\label{lem:mk-bisect}
Let $0<r<n$, $b_{ij}>0$, $d_j>0$. Starting from $\lambda_L=0$, $\lambda_H=D$ with the explicit
equilibria above, repeatedly query any box equilibrium on $[0,1]^n$ at the midpoint cutoff; return it
if its total is $r$, make the midpoint the new $\lambda_L$ if its total exceeds $r$, and the new
$\lambda_H$ if it is below $r$; record the pattern of the current endpoint equilibria. After at most
$1+\lceil\log_2(DQ_1^2)\rceil$ queries, either a BNE has been returned or $\lambda_H-\lambda_L<1/Q_1^2$.
In the latter case one computes by linear programming a rational $\lambda_0\in[0,D]$ and box equilibria
$(\Xlo{},f^L)$, $(\Xhi{},f^H)$ at capacities $c=\lambda_0/d$ with
\[
\Xlo{}\le \Xhi{},\qquad \sum_j\Xlo_j\le r\le\sum_j\Xhi_j ,
\]
all of polynomial binary length. All queries are at positive rational cutoffs of polynomial length.
\end{lemma}

\paragraph{Proof roadmap.} (1) Each endpoint's pattern polytope projects to an interval of cutoffs; (2) by
Cramer's rule these interval ends have denominators at most $Q_1$, so distinct ends differ by at least
$1/Q_1^2$; (3) once $\lambda_H-\lambda_L<1/Q_1^2$ the two intervals meet, and a common cutoff
$\lambda_0$ with two bracketing equilibria is found by LP; (4) Proposition~\ref{prop:mk-lattice}(b)
orders them.

\begin{proof}
A returned answer is a BNE by Lemma~\ref{lem:mk-box-cutoff}. Let $\Pi_L$ be the pattern polytope of the
lower endpoint with sign $\ge$ and $\Pi_H$ that of the upper endpoint with sign $\le$. Each contains its
endpoint equilibrium, so its projection onto $\lambda$ is a nonempty closed interval $I_L\ni\lambda_L$,
respectively $I_H\ni\lambda_H$, whose endpoints are computable by LP. Each endpoint of such an interval
is attained at a vertex of a pattern polytope, so by Cramer's rule and the Hadamard bound it is a
rational with denominator at most $Q_1$; two distinct such numbers, even from different patterns,
differ by at least $1/Q_1^2$.

When $\lambda_H-\lambda_L<1/Q_1^2$, the intervals $I_L$ and $I_H$ intersect. Otherwise, since
$\lambda_L<\lambda_H$, $I_L$ lies strictly to the left of $I_H$, and the right end of $I_L$ and the left
end of $I_H$ would be distinct such rationals in $[\lambda_L,\lambda_H]$, closer than $1/Q_1^2$. Choose
$\lambda_0$ in the intersection, for example the larger of the two left ends, and compute rational
points of $\Pi_L$ and $\Pi_H$ with $\lambda=\lambda_0$ by LP. These are box equilibria at $\lambda_0$
with totals $\ge r$ and $\le r$. Their coordinatewise minimum $\Xlo{}$ and maximum $\Xhi{}$ are box equilibria by
Proposition~\ref{prop:mk-lattice}(b), with $\sum \Xlo_j\le r\le\sum \Xhi_j$. Their bit lengths are polynomial
by the rational LP bounds. The number of halvings needed to bring $D$ below $1/Q_1^2$ is at most
$1+\lceil\log_2(DQ_1^2)\rceil$, and every midpoint is a positive rational number of polynomial length.
\end{proof}

The argument allows the oracle's selection to jump arbitrarily as the cutoff varies; an answer is used
only through its pattern. If an oracle returns a richer representation of an equilibrium, recovering
its positive support and price labels in polynomial time suffices.

\subsection{Narrowing and exact budget repair}

For a pattern as above, the \emph{price LP} has variables $x\in\mathbb R^n$ and $s$: keep the labels as
closed bounds ($x_j=0$, $x_j=1$, or $0\le x_j\le1$), the inequalities $U_{ij}(x)\ge U_{ik}(x)$ for all
allowed edges $(i,j)\in A$ and all $k$, add $0\le s\le n$ and $|\sum_jx_j-r|\le s$, and minimize $s$.
Its rows do not involve $\lambda$ or $f$. Bound the integer entries of all possible rows, after
clearing denominators, by $H_2\ge1$; with $N_2=n+1$ put $Q_2=(N_2H_2)^{N_2}$. A positive optimum is at
least $1/Q_2$. Put $\eta=1/(4nQ_2)$.

\emph{Canonicalization.} Given an equilibrium of a pinned sub-box, keep its positive support and its
actual labels on coordinates of positive width; on each zero-width coordinate fix the price and leave
the load unrestricted. Choose a vertex of the resulting price polytope (bounds, endpoint equalities,
and two-coordinate rows $U_{ij}(x)\ge U_{ik}(x)$) by rational LP, and independently a vertex of the
polytope of flows with this support, full rows, and the column conditions of the labels. If a formerly
free price reaches an endpoint of its interval, its load equals its capacity, which satisfies the new
label. The result still solves the pinned sub-box. In the procedure below, an answer is first
canonicalized and then patched by Lemma~\ref{lem:mk-pin}; the size bound of Lemma~\ref{lem:mk-bits}
needs this order.

\begin{lemma}[Narrowing and exact budget repair]\label{lem:mk-narrow}
In the situation of Lemma~\ref{lem:mk-bisect}, process coordinates in rounds. While some
$\Xhi_k-\Xlo_k>\eta$, choose a dyadic $t$ in the middle third of $[\Xlo_k,\Xhi_k]$ from a grid of spacing at most
$\eta/4$, query any equilibrium of the box market on $[\Xlo{},\Xhi{}]$ with $k$ pinned to $t$ (if $\lambda_0=0$,
use instead the upper corner of the pinned sub-box with any best responses), canonicalize it, and
recover an equilibrium $(x,f)$ of the box market on $[0,1]^n$ by Lemma~\ref{lem:mk-pin}. If
$\sum_jx_j=r$, return it; if $\sum_jx_j<r$ replace $\Xlo{}$ by $x$, otherwise replace $\Xhi{}$ by $x$. After
$O(n\log(1/\eta))$ queries every width is at most $\eta$. Then the price LP of the pattern of $\Xlo{}$ (with
its flow and the cutoff $\lambda_0$ kept fixed) has optimum zero, and any optimal $x'$ together with the
unchanged flow $f^L$ is an exact BNE.
\end{lemma}

\paragraph{Proof roadmap.} (1) Lemma~\ref{lem:mk-pin} keeps the invariants $\Xlo\le\Xhi$ and
$\sum\Xlo_j\le r\le\sum\Xhi_j$; (2) each query shrinks one width by a factor $2/3$; (3) once all
widths are at most $\eta$, $\Xlo$ is within $1/Q_2$ of the budget; (4) the Cramer gap of the price LP
then forces optimum zero; (5) an optimal $x'$ with the flow $f^L$ is a BNE by
Lemma~\ref{lem:mk-box-cutoff}.

\begin{proof}
At $\lambda_0=0$ all capacities vanish, so the upper corner is a box equilibrium of every sub-box. The
recovered $x$ lies in $[\Xlo{},\Xhi{}]$ and is an equilibrium on $[0,1]^n$, so the replacement keeps $\Xlo{}\le \Xhi{}$ and
$\sum \Xlo_j\le r\le\sum \Xhi_j$; a returned $x$ is a BNE by Lemma~\ref{lem:mk-box-cutoff}. Since $x_k=t$ is
in the middle third, the processed width shrinks by a factor at most $2/3$, and no width grows. A grid
point exists because the middle third has length more than $\eta/3$, and the grid needs
$O(\log(1/\eta))$ bits. So $O(n\log(1/\eta))$ queries make all widths at most $\eta$, and then
$|\sum_j\Xlo_j-r|\le\sum_j(\Xhi_j-\Xlo_j)\le n\eta<1/Q_2$. The point $\Xlo{}$ is feasible for the price LP of its own
pattern with $s<1/Q_2$, so the optimum is $0$. Let $x'$ be optimal. With the unchanged flow $f^L$ and
cutoff $\lambda_0$, every used edge is a best response at $x'$; zero and one labels keep their load
inequalities, and free labels have load exactly $c_j$, which is admissible at any price in $[0,1]$. So
$(x',f^L)$ is a box equilibrium at $\lambda_0$ with $\sum_jx'_j=r$, a BNE by
Lemma~\ref{lem:mk-box-cutoff}.
\end{proof}

\begin{lemma}[Sizes of adaptive queries]\label{lem:mk-bits}
With canonicalization, every endpoint and every query in Lemma~\ref{lem:mk-narrow} has binary length
polynomial in the input.
\end{lemma}

\paragraph{Proof roadmap.} A canonical price vertex is determined along spanning trees of tight comparisons
from one root. So each new coordinate is an affine image, with original-data coefficients, of a single
previous number (an endpoint coordinate or the pin), or of original data only if its component has no
active bound. Sizes therefore grow additively per call, not
multiplicatively.

\begin{proof}
Feeding general LP vertex bounds into the next LP could multiply the bit length at every step; the
canonical vertex avoids this. Make a graph on the coordinates of a canonical price vertex, joining two
distinct goods by a tight comparison $U_{ij}(x)=U_{ik}(x)$. In each connected component traverse a
spanning tree from a root. Each coordinate is an affine function of the root, a composition of at most
$n-1$ maps
\[
x_j=\frac{a_{ij}-a_{ik}}{b_{ij}}+\frac{b_{ik}}{b_{ij}}\,x_k
\]
built from original data only; multiplying all edge denominators along a path gives a common
denominator, and a sum of $n$ terms needs at most $\lceil\log_2(n+1)\rceil$ further bits.

If the
component has an active bound, take that coordinate as root: every value then involves only \emph{one}
previous endpoint coordinate or the dyadic pin. If it has none, some remaining tight comparison must
determine the root uniquely, since otherwise a small move of the whole component in both directions
preserves all active rows, contradicting that the point is a vertex; that equation uses original data
only.

Measuring the size of a rational by the larger bit length of its numerator and denominator, a map
$u\mapsto A+Bu$ with original-data coefficients increases the size additively by a fixed polynomial in
the input size. Hence the maximum endpoint size grows additively by a fixed polynomial per call, and
with polynomially many calls and polynomial initial sizes every query is polynomial. Flow vertices
depend only on the masses and capacities, and patching copies existing entries.
\end{proof}

\subsection{Zero slopes and zero gains}\label{sec:mk-degenerate}

\begin{lemma}[Degenerate data]\label{lem:mk-degenerate}
Let a game have $b_{ij}\ge0$, $d_j\ge0$, all $p_i>0$ and $0<r<n$. There is an explicitly computable
rational $\epsilon>0$ of polynomial bit length such that from any BNE of the game with
$b'_{ij}=b_{ij}+\epsilon$ and $d'_j=d_j+\epsilon$ one computes an exact BNE of the original game by one LP.
\end{lemma}

\paragraph{Proof roadmap.} (1) Relax each pattern LP by one slack $t$ and get a Cramer gap $1/Q_3$ for
positive optima; (2) set $\epsilon=1/(8Q_3)$; (3) a BNE of the perturbed game satisfies its own pattern
at the original data with slack at most $\epsilon$; (4) hence that pattern's slack is $0$, and an optimal
solution is an exact BNE.

\begin{proof}
For a pattern (allowed edges and zero/one/free labels), use the original variables $(x,f,\lambda)$ of
the pattern polytope with the \emph{exact} budget $\sum_jx_j=r$. Introduce a slack $t$ with
$0\le t\le1$, split each free-target equality into two inequalities, and relax every pairwise
best-response inequality and every target cutoff inequality by $t$, in unscaled payoff units. Keep
exact all flow bounds, row sums, zeros off allowed edges, price bounds, label equalities and the
budget; bound $0\le\lambda\le D$ with $D=1+\max_jd_j$. After inserting the slack coefficients, clear
denominators row by row over the full menu of possible rows, let $H_3\ge1$ bound all integer entries,
and with $N_3=n+mn+2$ put $Q_3=(N_3H_3)^{N_3}$. The LP is bounded, so by Cramer's rule every positive
minimum of $t$ over every pattern is at least $1/Q_3$. Set $\epsilon=1/(8Q_3)$.

Take a BNE of the perturbed game and its actual support and labels. At the original payoffs, the
pairwise difference $U_{ij}-U_{ik}$ changes by $\epsilon(x_j-x_k)$, so a used good falls short by at
most $\epsilon$. Each score changes by $\epsilon w_j\le\epsilon$. A positively covered and an
unsaturated target exist, and every perturbed score lies in $[0,\max_jd'_j]\subseteq[0,D]$, so the
perturbed equilibrium has a cutoff in $[0,D]$. The original pattern LP is therefore feasible with
$t=\epsilon<1/Q_3$, its minimum is $0$, and a rational optimal solution is an exact BNE of the original
game. Formerly positive edges that lose their flow remain allowed tight edges, which preserves all
best-response conditions.
\end{proof}

\subsection{The game-to-market direction}

\begin{proposition}[Game to market]\label{prop:mk-g2m}
\statusrefereed
Theorem~\ref{thm:mk-main}(b) holds: polynomially many calls to an oracle returning any equilibrium of
a dense (M1) instance, each call of polynomial binary length and each answer assumed of polynomial
length, together with polynomial-time rational linear programming, compute an exact BNE of any explicit
game.
\end{proposition}

\begin{proof}
Drop zero-prior types and handle $r\in\{0,n\}$ directly. If some slope or gain is zero, pass to the
perturbed game of Lemma~\ref{lem:mk-degenerate}, whose data have polynomial length. For the positive
game, run Lemma~\ref{lem:mk-bisect} and Lemma~\ref{lem:mk-narrow}. Every box query is made at a positive
cutoff (bisection midpoints are positive and at $\lambda_0=0$ no query is made), so all capacities are
positive and Lemma~\ref{lem:mk-box-m1} turns it into one dense (M1) call in zero-reserve form whose
every answer decodes. The number of calls is $O(\log(DQ_1^2)+n\log(1/\eta))$, which is polynomial, and
each call has polynomial length by Lemma~\ref{lem:mk-bits}. Decode by Lemma~\ref{lem:mk-degenerate} if
needed, realize the coverage by Lemma~\ref{lem:realize}, and give zero-prior types any best response.
\end{proof}

The (M1) calls made here have the complete edge set and arbitrary positive slopes; they do not stay
within a restricted slope class. The reduction is adaptive, so it is a Turing reduction, not a many-one
reduction.

\section{Market to game}\label{sec:mk-m2g}

\paragraph{Idea.} Each good becomes a target and each buyer a type whose payoff at target $j$ is its
market payoff, scaled down by a constant $S$ so that coverages stay small. Three kinds of extra types
force the defender to behave like a price-setter. A \emph{dedicated} type per good puts weight
$W-c_j$ on that target, so a target is ``full'' (load $W$) exactly when buyers fill its capacity $c_j$. A
\emph{control} type guarantees that the maximal load is exactly $W$. A \emph{private} target per buyer
absorbs all of the buyer's weight beyond its mass $\mu_i$ and encodes its utility $v_i$ in its
coverage. One condition is essential: all slopes must be at least one.

\subsection{Reducing to slopes at least one}

\begin{lemma}[Row rescaling]\label{lem:mk-rescale}
\statuslean
Given an (M1) instance with positive slopes, multiply all intercepts and slopes of each buyer $i$ by a
positive power of two $\nu_i$ with $\nu_i\min_{j\in E_i}b_{ij}\ge1$. Then $(P,v,f)$ is an equilibrium of
the original instance if and only if $(P,(\nu_iv_i)_i,f)$ is an equilibrium of the rescaled one. Each
$\nu_i$ has polynomial binary length.
\end{lemma}

\begin{proof}
Conditions (a) and (b) for buyer $i$ are multiplied by $\nu_i>0$, and $v_i>0$ if and only if
$\nu_iv_i>0$; conditions (c) and (d) involve only $f$, $P$ and the sign of $v_i$. The exponent of
$\nu_i$ is bounded by the bit length of $\min_jb_{ij}$.
\end{proof}

\subsection{The one-resource embedding}

Let an (M1) instance have buyers $I$, goods $J$, and slopes $b_{ij}\ge1$ on all
listed pairs. Fix a rational $0<\omega\le1$. Put
\begin{equation}\label{eq:mk-consts}
\hat D=1+\max\Bigl(\sum_i\mu_i,\ \max_jc_j\Bigr),\quad W=\Bigl(2+\frac2\omega\Bigr)\hat D,\quad
A=1+\max\bigl(0,\max_{ij}a_{ij}\bigr),\quad S=1+(|I|+|J|)A .
\end{equation}
The game $\Gamma_\omega$ has a \emph{real} target $j$ for each good, a \emph{private} target $g_i$ for
each buyer, one \emph{outside} target $o$ and one \emph{control} target $h$; one defender resource; and
defender payoffs $0$ at an uncovered and $1$ at a covered target, for every target. Write $\xi_k$ for
the coverage of target $k$. Every attacker payoff below has the form $u_2^0-\beta\xi_k$ with slope
$\beta\ge0$, i.e.\ $u_2^1=u_2^0-\beta$. The types, with unnormalized weights, are listed in
Table~\ref{tab:mk-types}.
\begin{table}[h]
\centering
\begin{tabular}{lll}
\toprule
type & weight & payoffs \\
\midrule
source $i\in I$ & $W+\mu_i$ & $a_{ij}/S-b_{ij}\xi_j$ at real $j\in E_i$; $A/S-\xi_{g_i}$ at $g_i$; $-\xi_o$ at $o$\\
dedicated $j\in J$ & $W-c_j$ & $2-\xi_j$ at real $j$\\
control & $W$ & $2-\xi_h$ at $h$\\
\bottomrule
\end{tabular}
\caption{The types of $\Gamma_\omega$. Every other (type, target) pair pays $-2-\xi_k$.}
\label{tab:mk-types}
\end{table}

In particular a non-listed pair $(i,j)$, $j\notin E_i$, pays $-2-\xi_j$. The prior is the weight vector
divided by its total. There are $|I|+|J|+1$ types and $|I|+|J|+2$ targets. By the one-resource remark of
Section~\ref{sec:pre-game}, a BNE is exactly a profile in which every type mixes over its maximizing
targets and $\supp\xi\subseteq\argmax_k\ell_k$, where $\ell_k$ is the total weight attacking $k$.

\begin{theorem}[One-resource embedding]\label{thm:mk-embed}
\statusrefereed
Let the (M1) instance have all listed slopes $b_{ij}\ge1$. Then $\Gamma_\omega$ is a legal Bayesian
simple security game, computable in time polynomial in the market and the length of $\omega$, with
one resource, defender payoffs $0$ and $1$, attacker slopes equal to $b_{ij}$ on listed source$\times$real
pairs and $1$ elsewhere, and priors with $1-\omega\le|\Theta|\Pr(\theta)\le1+\omega$ (where $|\Theta|$ is its number of types) and
largest-to-smallest ratio at most $1+\omega$. For \emph{every} BNE of $\Gamma_\omega$, let $f_{ij}$
be the unnormalized weight that source type $i$ puts on real target $j$ and $u_i$ that type's
equilibrium payoff. Then
\[
P_j=S\xi_j,\qquad v_i=Su_i,\qquad f
\]
form an equilibrium of the market. Conversely every market equilibrium lifts to a BNE, and the two maps
are mutually inverse between market price vectors and BNE coverage vectors.
\end{theorem}

\paragraph{Proof roadmap.} (1) Legality and the prior bounds are arithmetic. (2) Decoding has four steps:
the control type pins the maximal load to at least $W$; private targets absorb at least $W$ of each
source; the coverage bound (the only use of $b_{ij}\ge1$) shows the control target is covered, so the
maximal load is exactly $W$; then conditions (a)--(d) are read off. (3) The lift constructs a BNE from a
market equilibrium. (4) The two maps are inverse.

\begin{proof}
\emph{Legality and priors.} All slopes are nonnegative and $0\le1$ for the defender. Every weight lies
in $[W-\hat D,W+\hat D]$, and $(W+\hat D)/(W-\hat D)=1+2\omega/(\omega+2)\le1+\omega$; the average
weight also lies in this interval, which gives both prior bounds, using $(1+\omega)^{-1}\ge1-\omega$.
All numbers have polynomial length.

\emph{Decoding.}
\begin{enumerate}[nosep,label=\arabic*.]
\item \emph{Control.} The control type gets $2-\xi_h\ge1$ at $h$ and at most $-2$ elsewhere, so it
attacks $h$. No other type attacks $h$: a source type has the outside target worth
$-\xi_o\ge-1>-2-\xi_h$, and a dedicated type has its own target worth at least $1$. So $\ell_h=W$, and
every covered target has load $\max_k\ell_k\ge W$. Likewise a dedicated type attacks only its own
target, only source $i$ attacks $g_i$, and no non-listed pair is ever used.
\item \emph{Private targets.} If $\xi_{g_i}>0$, then $g_i$ has load at least $W$, all from source $i$.
If $\xi_{g_i}=0$, then $g_i$ pays $A/S$, which strictly beats every real option
($a_{ij}/S-b_{ij}\xi_j\le a_{ij}/S<A/S$; this uses only $b_{ij}\ge0$), the outside ($\le0$) and every
unlisted target ($\le-2$); so source $i$ puts all $W+\mu_i$ on $g_i$. In both cases at least $W$ goes to
$g_i$. The source mass left for real targets and $o$ is at most $\sum_i\mu_i<W$. Only source types
attack $o$, so $\ell_o<W$ and $\xi_o=0$. Hence every source type has $u_i\ge0$.
\item \emph{Coverage bound.} If $\xi_j>0$ for a real $j$, then $\ell_j\ge W$ while the dedicated weight
is $W-c_j$, so some source $i$ attacks $j$; from $a_{ij}/S-b_{ij}\xi_j=u_i\ge0$ we get
$\xi_j\le a_{ij}/(Sb_{ij})<A/S$. \emph{This is the only step of the decoding that uses $b_{ij}\ge1$.}
If $\xi_{g_i}>0$ its owner attacks it, so $\xi_{g_i}\le A/S$. Hence
$\sum_j\xi_j+\sum_i\xi_{g_i}\le(|I|+|J|)A/S<1$; with $\xi_o=0$ this gives $\xi_h>0$, so
$\max_k\ell_k=\ell_h=W$. An uncovered private target would carry $W+\mu_i>W$, so every $\xi_{g_i}>0$,
every $g_i$ receives exactly $W$, and source $i$ places exactly $\mu_i$ on real targets and $o$.
\item \emph{Reading off.} For a listed pair, $u_i\ge a_{ij}/S-b_{ij}\xi_j$ is condition (a) after
multiplying by $S$, since $a_{ij}/S-b_{ij}\xi_j=(a_{ij}-b_{ij}P_j)/S$; tightness of used pairs is (b).
If $v_i>0$ then $u_i>0=-\xi_o$, so nothing goes to $o$ and $\sum_jf_{ij}=\mu_i$, which is (c). Finally
$\sum_if_{ij}=\ell_j-(W-c_j)\le c_j$, with equality when $\xi_j>0$, which is (d).
\end{enumerate}

\emph{Lift.} Given a market equilibrium, $v_i=\max(0,\max_{j\in E_i}(a_{ij}-b_{ij}P_j))$ by (a)--(c). A
good with $P_j>0$ is full and has a user with nonnegative payoff, so $0\le P_j\le a_{ij}/b_{ij}<A$ (using
$b_{ij}\ge1$), and $0\le v_i<A$. Set
\[
\xi_j=P_j/S,\qquad \xi_{g_i}=(A-v_i)/S,\qquad \xi_o=0,\qquad \xi_h=1-\sum_j\xi_j-\sum_i\xi_{g_i}>0 .
\]
Give $W$ units of source $i$ to $g_i$ (worth $v_i/S$, its row maximum), the market flow $f$ to real
targets, and the rest to $o$ (worth $0$, used only if $v_i=0$ by (c)); dedicated and control types use
their own targets. Real loads are at most $W$, with equality when covered by (d); private and control
loads are $W$; the outside load is below $W$. So every covered target has maximum load and all types
best respond. Dividing flows by type weights gives conditional attack distributions.

\emph{Inverse maps.} In every BNE, step 3 shows $\xi_o=0$ and that source $i$ uses $g_i$, so
$u_i=A/S-\xi_{g_i}$, i.e.\ $\xi_{g_i}=(A-v_i)/S$, and $\xi_h$ is determined by the rest. Since $v$ is
determined by $P$, the coverage vector of a BNE is the lift of its decoded price vector, and decoding a
lift returns $P$.
\end{proof}

The proof uses $W$ only through $W>\sum_i\mu_i$ and $W>\max_jc_j$. For $\omega=1$ the base weight is
$W=4\hat D$; this is the case used in Chapter~\ref{ch:arrival}. Since $\Gamma_\omega$ is finite it
has a BNE, so every (M1) market with slopes $\ge1$ has an equilibrium; with
Lemma~\ref{lem:mk-rescale} this extends to all positive slopes.

\begin{proposition}[The embedding fails for a slope below one]\label{prop:mk-below-one}
\statusrefereed
Consider the market with one buyer of mass $2$, one good of capacity $1$, and payoff $1-P/100$ (slope
$1/100$). Its unique equilibrium has $P=100$. Apply the construction literally with $\omega=1$, so
$\hat D=3$, $W=12$, $A=2$, $S=5$, and order the targets as (real, private, outside, control). Then
\[
\xi=\Bigl(\tfrac{80}{101},\ \tfrac{21}{101},\ 0,\ 0\Bigr)
\]
is a BNE with the control uncovered, and it decodes to $P=S\xi_1=400/101$, which is not a market
equilibrium.
\end{proposition}

\begin{proof}
For $P<100$ demand $2$ exceeds capacity $1$, and for $P>100$ the buyer's payoff is negative, so
$P=100$ is the only equilibrium price. In the game, the source type (weight $14$) gets
$1/5-\xi_1/100=97/505$ at the real target, $2/5-\xi_2=97/505$ at its private target, $0$ at the outside
and at most $-2$ elsewhere, so it is tied between real and private. Send $3/2$ of its
weight to the real target and $25/2$ to the private one. The dedicated type (weight $11$) attacks the
real target and the control type (weight $12$) the control. Loads are $(25/2,25/2,0,12)$: both covered
targets carry the maximum load $25/2$, so the defender best responds. The decoded price $400/101$ is not
$100$.
\end{proof}

So the hypothesis $b_{ij}\ge1$ cannot be dropped from Theorem~\ref{thm:mk-embed}. We handle arbitrary
positive slopes by the row rescaling of Lemma~\ref{lem:mk-rescale}.

\section{Proof of the main theorem and the corollary}\label{sec:mk-proof}

\begin{proof}[Proof of Theorem~\ref{thm:mk-main}]
(b) is Proposition~\ref{prop:mk-g2m}. For (a), given an (M1) instance with arbitrary positive slopes,
rescale rows by Lemma~\ref{lem:mk-rescale} so that all slopes are at least one, build
$\Gamma_\omega$ (say with $\omega=1$) by Theorem~\ref{thm:mk-embed}, call the BNE oracle once,
decode, and undo the rescaling of the utilities. All steps take polynomial time.
\end{proof}

\begin{proof}[Proof of Corollary~\ref{cor:mk-one-resource}]
Given a general game, run the algorithm of Proposition~\ref{prop:mk-g2m}. Answer each of its
polynomially many (M1) calls by rescaling rows (Lemma~\ref{lem:mk-rescale}), building $\Gamma_\omega$
(Theorem~\ref{thm:mk-embed}), calling the restricted BNE oracle once, and decoding. Each
$\Gamma_\omega$ has one resource, defender payoffs $0$ and $1$ and the stated priors, and its length is
polynomial in the call and in the length of $\omega$. The reverse reduction is inclusion.
\end{proof}

\section{What the equivalence says}

The equivalence is a reduction in both directions. By itself it is neither a polynomial algorithm nor
a hardness result: it isolates the complexity of exact equilibrium search in general (M1) markets as the
question. Hardness proved for either formulation transfers to the other. It concerns \emph{exact}
equilibria: the decoding uses the exact equilibrium conditions, and the bisection and budget repair use
exact rational separation bounds. A balanced matching market with a prescribed price anchor (every row
and column filled, signed prices, one price fixed at zero) is a \emph{different} model, into which (M1)
with zero reserves embeds; neither direction of Theorem~\ref{thm:mk-main} uses it.
